%!TEX root = ../template.tex
%%%%%%%%%%%%%%%%%%%%%%%%%%%%%%%%%%%%%%%%%%%%%%%%%%%%%%%%%%%%%%%%%%%%
%% abstrac-pt.tex
%% ISEL thesis document file
%%
%% Abstract in Portuguese
%%%%%%%%%%%%%%%%%%%%%%%%%%%%%%%%%%%%%%%%%%%%%%%%%%%%%%%%%%%%%%%%%%%%
Um \emph{workflow} \emph{serverless} é escrito antes de existirem as funções que invoca. O
\emph{endpoint} de uma função --- um URL do Cloud Run, um ARN do Lambda --- é produzido pela sua
implantação e depende da conta, do projeto, da região e do fornecedor onde foi implantada. O
programador implanta as funções com uma ferramenta, lê os \emph{endpoints} produzidos e copia-os
à mão para a definição do \emph{workflow}, consumida por outra. A cópia repete-se sempre que a
função muda de região, de conta ou de fornecedor, e um \emph{endpoint} mal copiado só falha em
execução.

Esta dissertação elimina esse passo unificando as duas ferramentas: o QuickFaaS, que implanta
funções \emph{serverless} portáveis, e o OmniFlow, uma DSL em Kotlin que gera e implanta
\emph{workflows} no AWS Step Functions e no Google Cloud Workflows. Uma chamada identifica a
função por um identificador lógico estável, e uma cascata de resolução em três níveis liga-o no
momento da implantação: consulta um registo de
funções por fornecedor e valida a entrada contra a função existente; se falhar, procura a função
nas regiões do fornecedor; e só com a ausência confirmada a implanta através do QuickFaaS, a partir
de um descritor. O QuickFaaS, até aqui limitado ao Google Cloud e ao Azure, foi estendido com um
fornecedor para o AWS Lambda, cobrindo as mesmas duas \emph{clouds} que o OmniFlow.

O protótipo é avaliado quanto à correção e ao custo. Quarenta e cinco testes unitários
cobrem todos os caminhos da cascata nos dois fornecedores, e os testes de desempenho mostram que a
resolução local custa milissegundos em \emph{workflows} com até duzentas chamadas internas,
crescendo linearmente com esse número e não com a dimensão do registo. Implantar
uma aplicação \emph{serverless} a partir de uma única definição é, assim, exequível, e a resolução
acrescentada é pequena face à implantação na \emph{cloud}.

% Palavras-chave do resumo em Português
\begin{keywords}
Computação \emph{serverless}, \emph{Function-as-a-Service}, Orquestração de \emph{workflows},
Portabilidade \emph{multi-cloud}, Ligação de \emph{endpoints} em tempo de implantação, AWS Lambda,
Google Cloud
\end{keywords}
